%! Characteristics of Functions: Domain, Range, Intercepts, and End Behavior — Unit 2, Lesson 2.4 — Warm-Up
% Exactly ONE page, blank and keyed (LESSON_SHAPE.md, one_page). Five prerequisite items, each
% the algebra one row of today's notes runs on: 1–2 solve f(x) = 0 by factoring and by square
% roots (INTERCEPTS), 3 solves the inequality a radicand must satisfy (FROM A FORMULA), 4 is the
% interval notation 1.1 met lightly (every row), and 5 is the sign of a big input's output
% (END BEHAVIOR).
\documentclass[10pt]{article}
\usepackage{precalculus-article}
\usepackage{precalculus-boxes}
\newcommand{\TallMath}[1]{$\displaystyle #1\rule[-1.4em]{0pt}{3.2em}$}

\begin{document}

\pageheader{Unit 2, Lesson 2.4}{Warm-Up}

\vspace{0.12in}

\begin{enumerate}[itemsep=10pt]

\item Solve by factoring: \quad $x^2 - 5x - 6 = 0$

\begin{work}
  x^2 - 5x - 6 &= 0 \\
  (x - 6)(x + 1) &= 0 \\
  x = 6 \quad\text{or}\quad x &= -1
\end{work}

$x = $ \blank{3.0cm}

\item Solve by taking square roots: \quad $x^2 - 49 = 0$

\begin{work}
  x^2 &= 49 \\
    x &= \pm 7
\end{work}

$x = $ \blank{3.0cm}

\item Solve the inequality: \quad $3x - 12 \ge 0$

\begin{work}
  3x - 12 &\ge 0 \\
       3x &\ge 12 \\
        x &\ge 4
\end{work}

Solution: \blank{3.0cm}

\item Write each set of numbers in interval notation. A square bracket $[\ ]$ includes the
endpoint; a round parenthesis $(\ )$ leaves it out.

\begin{enumerate}[label=(\alph*), itemsep=8pt]
  \item $-2 < x \le 5$ \qquad \blank{2.6cm}
  \item $x \ge 4$ \qquad \blank{2.6cm}
\end{enumerate}

\item Let $f(x) = -x^2$. Find $f(10)$ and $f(-10)$. Is each output positive or negative?

\begin{work}
  f(10) &= -(10)^2 = -100 \\
  f(-10) &= -(-10)^2 = -100
\end{work}

$f(10) = $ \blank{2.0cm} \qquad $f(-10) = $ \blank{2.0cm} \qquad Sign: \blank{2.4cm}

\end{enumerate}

\end{document}
